\begin{Example}
\label{ex:electrostatic-surface-charge-weak-solution}
An important series of classical examples of Dirichlet problems come from electrostatics. Without going into any of the physics, let us assume that we wish to know a scalar quantity $u$ called the \emph{electrostatic potential} or more commonly the \emph{voltage} in a domain $\Omega\subset\mathbb{R}^3$. We will assume that the boundary of the domain is grounded; i.e.,
\begin{equation}\label{eq:electrostatic-grounded-dirichlet-boundary}
u(\bm{x})=0\quad\text{for }\bm{x}\in\partial\Omega.
\end{equation}
Within $\Omega$ there is a distribution of charge $\rho:\Omega\to\mathbb{R}$, and this charge ``generates\textquotesingle\textquotesingle\ a voltage through the formula
\begin{equation}\label{eq:electrostatic-charge-poisson-equation}
-\Delta u=\rho.
\end{equation}
The solution of this class of problems is the subject of classical potential theory, and many of the techniques of the classical theory (eigenfunction expansion, Green\textquotesingle s functions) are included in the modern theory as well. However, the classical and modern theory share the same approach only when we are looking for classical ($\rho\in C_b(\Omega)$) or strong ($\rho\in L^2(\Omega)$ or piecewise continuous) solutions.

The classical and modern theories take a very different approach in dealing with charge distributions that occur on surfaces; i.e., situations where $S\subset\Omega$ is a smooth surface and a charge distribution $\omega:S\to\mathbb{R}$ is defined. (We assume that either $\omega\in C_b(S)$ or $\omega\in L^2(S)$.) Of course, since $\omega$ is defined only on a surface (a set of measure zero to readers who have had measure theory) the differential equation $-\Delta u=\omega$ does not make sense either classically or as the identification of equivalence classes of sequences in $L^2(\Omega)$. In the modern theory the situation is very clear: although $\omega$ is not in $L^2(\Omega)$, we can use it to define a perfectly nice functional in $H^{-1}(\Omega)$ through the formula
\begin{equation}\label{eq:surface-charge-negative-sobolev-functional}
(\omega,\phi)=\int_S\omega(\bm{x})\phi(\bm{x})\,\mathrm{d}a(\bm{x}),
\end{equation}
for every $\phi\in H_0^1(\Omega)$. (Recall that the trace theorem implies that $\phi\in L^2(S)$.) Here $\mathrm{d}a(\bm{x})$ indicates differential area at $\bm{x}\in S$. Thus, the modern theory would simply have us look for a (weak) solution $u\in H_0^1(\Omega)$ of the variational problem
\begin{equation}\label{eq:surface-charge-electrostatic-variational-problem}
B[v,u]:=\int_\Omega\nabla v\cdot\nabla u\,\mathrm{d}\bm{x}=(\omega,v)
\end{equation}
for all $v\in H_0^1(\Omega)$. As we shall see below, this problem is well-posed. The classical theory solves this problem (and some other similar ones) using the theory of single and double layer potentials: essentially integral operators defined using singular surface integrals. As is so often the case, the classical theory lacks much of the conceptual unity of the modern theory, but provides much more detailed information in special (though often the most important) cases. We will not go into the results of classical potential theory in this book, but the reader is encouraged to read \emph{Foundations of Potential Theory} by O.D. Kellogg \cite{Ke} as a good starting point for more information on this subject.
\end{Example}
